% !TEX root = ./Basilisk-avsLibrary20170812.tex

\section{Test Description and Success Criteria}
The AVS support libraries undergo a series of unit tests within C-functions that are called by the Swig'd interface {\tt avsLibrarySelfCheck}.  The automated unit test is performed via {\tt pytest} through the script:
\begin{verbatim}
     architecture/utilitiesSelfCheck/_UnitTest/test_avsLibrarySelfCheck.py
\end{verbatim}
In each unit test the C-code library function output is compared to an externally hand-computed solution.  The success criterion is when the pre-computed solution and the C-code computed solution match to within a certain numerical tolerance.  The success criteria absolute value tolerances are shown in Table~\ref{tbl:tolerance}.



\begin{table}[htbp]
	\caption{Unit Test Absolute Difference Tolerances}
	\label{tbl:tolerance}
	\centering \fontsize{10}{10}\selectfont
	\begin{tabular}{c | c} % Column formatting,
		\hline
		\hline
		Test    & Absolute Tolerance \\
		\hline
		  Linear Algebra    & $10^{-10}$ \\
		  Rigid Body Kinematics    & $10^{-10}$ \\
		  Orbital Motion: Anomalies    & $10^{-10}$ \\
		  Orbital Motion: Hill transformations    & $10^{-4}$ \\
		\hline
		\hline
	\end{tabular}
\end{table}





\section{Test Parameters}

\subsection{Linear Algebra}
In these unit tests each function is provided non-trivial inputs, such as {\tt v1 = [1,2,3]} to evaluate the associate matrix math.  For specific information on the values used, look at the source code for the function {\tt testLinearAlgebra()} in the file {\tt avsLibrarySelfCheck.c}.

\subsection{Rigid Body Kinematics}
The attitude kinematics sub-routines are provided with non-trivial inputs to test the functions.   For specific information on the values used, look at the source code for the function {\tt testRigidBodyKinematics()} in the file {\tt avsLibrarySelfCheck.c}.


\subsection{Orbital Motion}
The orbital anomaly sub-routines are provided with non-trivial inputs to test the functions.   For specific information on the values used, look at the source code for the function {\tt testOrbitalAnomalies()} in the file {\tt avsLibrarySelfCheck.c}.

The Hill self-check, {\tt testOrbitalHill()}, exercises rotating-frame state
transformations. The native target {\tt test\_orbitalMotion} checks element
conversions and environment models. Additional Python regressions are under
{\tt src/tests}, with integration tests in module test directories. These are separate from the C self-check entry points. Individual
assertions define tolerances appropriate to each quantity and case.
Compatibility regressions include repeated calls with Python and SWIG element
containers, ignored auxiliary fields in the legacy parabolic adapter, and C
compound-literal and C++ braced-initializer zonal calls. Large wrapped anomalies
are checked against direct trigonometric state references in both languages.

\section{Test Results}
The four native self-check entry points exercised by this Python wrapper are listed in Table~\ref{tbl:avsLibraryResults}. Unknown selectors fail explicitly; no result is claimed for an entry point that does not exist.


\begin{table}[h]
	\caption{Integration test results.}
	\label{tbl:avsLibraryResults}
	\centering \fontsize{10}{10}\selectfont
	\begin{tabular}{c | c | l } % Column formatting,
		\hline\hline
		\textbf{Test} 			& \textbf{Pass/Fail} 	 & \textbf{BSK Error Notes}
		\\ \hline
		{\tt rigidBodyKinematics}		  	&
		\input{AutoTeX/test_avsLibrarySelfCheckTestMsg-testRigidBodyKinematics}      	  &
		\input{AutoTeX/test_avsLibrarySelfCheckMsg-testRigidBodyKinematics}
		\\
		{\tt orbitalHill}		  	&
		\input{AutoTeX/test_avsLibrarySelfCheckTestMsg-testOrbitalHill}      	  &
		\input{AutoTeX/test_avsLibrarySelfCheckMsg-testOrbitalHill}
		\\
		{\tt orbitalAnomalies}		  	&
		\input{AutoTeX/test_avsLibrarySelfCheckTestMsg-testOrbitalAnomalies}      	  &
		\input{AutoTeX/test_avsLibrarySelfCheckMsg-testOrbitalAnomalies}
		\\
		{\tt linearAlgebra}		  	&
		\input{AutoTeX/test_avsLibrarySelfCheckTestMsg-testLinearAlgebra}      	  &
		\input{AutoTeX/test_avsLibrarySelfCheckMsg-testLinearAlgebra}
		\\
		\hline\hline
	\end{tabular}
\end{table}
